%%%%%%%%%%%%%%%%%%%%%%%%%%%%%%%%%%%%%%%%%%%%%%%%%%%%%%%%%%%%%%%%%%%%%%
% Lớp 12 | Ứng dụng đạo hàm để khảo sát hàm số | Giá trị lớn nhất và giá trị nhỏ nhất của hàm số
% Dạng 2D1?3-5: GTLN, GTNN hàm nhiều biến
%%%%%%%%%%%%%%%%%%%%%%%%%%%%%%%%%%%%%%%%%%%%%%%%%%%%%%%%%%%%%%%%%%%%%%


% nguồn: 2026-11-16_12_2D1_TL_Ham_so (nhập 16/11/2026)
\begin{bt}%[2D1V3-5]%[ma:ec62962a]
% nguon: Archive/SP_KhoLatex/ID6_old/DeThamKhao-03.tex #3680
%[Dự án ĐGNL-ĐGTD-2022 GĐ1]%[Duong Xuan Loi]
	Cho $x$, $y>0$ thỏa mãn $\heva{& x^2-xy+3=0 \\ & 2x+3y-14\le 0}$. Tính tổng giá trị lớn nhất và nhỏ nhất của biểu thức $P=3x^2y-xy^2-2x^3+2x$.
	
	\textit{Kết quả:} \tikz\draw (0,0) rectangle (3,0.6);
	\loigiai{
		Ta có $x^2-xy+3=0$ $\Leftrightarrow y=\dfrac{x^2+3}{x}$.\\
		Khi đó $2x+3y-14\le 0 \Leftrightarrow 2x+\dfrac{3x^2+9}{x}-14\le 0 \Leftrightarrow 5x^2-14x+9\le 0 \Leftrightarrow x\in \left[1;\dfrac{9}{5}\right]$.\\
		Cũng có $P=3x^2\cdot \dfrac{x^2+3}{x}-x\cdot {\left(\dfrac{x^2+3}{x}\right)}^2-2x^3+2x=5x-\dfrac{9}{x}=f(x)$.\\
		Mà $f'(x)=5+\dfrac{9}{x^2}>0$, $\forall x\in \left[1;\dfrac{9}{5}\right]$.\\
		Do đó GTLN của $P$ là $f\left(\dfrac{9}{5}\right)=4$ và GTNN của $P$ là $f(1)=-4$.\\
		Vậy tổng giá trị lớn nhất và nhỏ nhất của biểu thức $P$ bằng $0$.
	}
\end{bt}
